\documentclass[10pt,aspectratio=169]{beamer}
\usepackage{graphicx}

%\ifxetex
\usepackage{fontspec}
\newfontfamily{\klavikumfont}{Klavika-Regular.otf}
\selectfont\klavikumfont
%\fi

\usepackage[yyyymmdd]{datetime}
\renewcommand{\dateseparator}{-}

\setbeamercolor{normal text}{fg=white,bg=black}
\setbeamerfont{author}{size=\small,series=\bfseries,parent=structure}
\setbeamerfont{frametitle}{family=\fontfamily{cyklop}\selectfont}
\setbeamercolor{background canvas}{bg=white}

\pgfdeclareimage[width=\paperwidth]{mybackground}{gpu_background_black.png}
\setbeamertemplate{title page}{
  \begin{picture}(0,0)
    \put(-30,-163){%
      \pgfuseimage{mybackground}
    }
    \put(0,-150){%
      \begin{minipage}[b][45mm][t]{226mm}
        \usebeamerfont{title}{\inserttitle\par}
        \ifx\insertsubtitle\@empty%
        \else%
        {\vskip0.25em%
        \usebeamerfont{subtitle}\usebeamercolor[fg]{subtitle}\insertsubtitle\par}%
        \fi%
        \vskip1cm%              
        \usebeamerfont{author}{\insertauthor}\par
        \usebeamerfont{date}\insertdate
      \end{minipage}
    }
  \end{picture}
}

\setbeamerfont{frametitle}{family=\fontspec{Klavika-Regular.otf}}

\def\thetitle{rocFFT}
\def\thesubtitle{Performance updates}
\def\theauthor{rocFFT developers}
\def\theauthorinstitute{AMD}

%\def\presdate{2020-02-14}

\hypersetup{
  pdftitle=\thetitle{} \thesubtitle{},
  pdfauthor=\theauthor{}
}

\title{\thetitle{}}
\subtitle{\thesubtitle{}}
\author{\theauthor{}}
\institute{\theauthorinstitute}
%\date{\presdate}

\setbeamertemplate{frametitle}[default][center]
\setbeamertemplate{navigation symbols}{}

\begin{document}

\klavikumfont
\setbeamertemplate{headline}{}
%\setbeamertemplate{footline}{}

{
  \setbeamercolor{background canvas}{bg=black}
  \setbeamercolor{normal text}{fg=white,bg=black}
  \setbeamercolor{frametitle}{fg=white}
  \setbeamercolor{title}{fg=white}
  \setbeamercolor{alerted text}{fg=white}
  \setbeamercolor{structure}{fg=white}
  \usebeamercolor[fg]{normal text}
  \begin{frame}
    \titlepage
  \end{frame}
}

\setbeamercolor{normal text}{fg=black,bg=white}
\usebeamercolor[fg]{normal text}

\section{Overview}

\begin{frame}
  \frametitle{rocFFT performance update}

  Overview
  \begin{itemize}
  \item Brief intro to FFTs
  \item Performance numbers 2021-06-29 \hyperlink{perf01}{\beamergotobutton{skip to}}
  \item Performance planning
  \end{itemize}
\end{frame}

\section{Intro to FFTs}

\begin{frame}
  \frametitle{Intro to FFTs}

  Given length $N$ input data $x_n$, discrete Fourier transform (DFT) is
  \begin{equation*}
    X_k = \sum_{n=0}^{N-1} e^{-2 \pi i n k / N} x_n
  \end{equation*}

  Note
  \begin{itemize}
  \item output is also length $N$
  \item each output element depends on all input elements
  \end{itemize}
\end{frame}

\begin{frame}[fragile]
  \frametitle{Naive DFT}

  A naive approach to compute the DFT would be

\begin{verbatim}
X = zeros(N)
for k in range(N):
  for n in range(N):
    X[k] += exp(-2 * pi * I * n * k / N) * x[n]
\end{verbatim}

  Properties
  \begin{itemize}
  \item slow: code above is $O(N^2)$
  \item not very accurate when done with floating-point
  \end{itemize}
\end{frame}

\begin{frame}
  \frametitle{Fast Fourier Transform}

  Key idea behind Fast Fourier Transform (FFT)
  \begin{itemize}
  \item lots of symmetries in the $e^{-2 \pi i n k / N}$ term
  \item symmetries manifest themselves through factors of $N$
  \end{itemize}

  \medskip
  For example, length 56 transform is decomposed into
  \begin{itemize}
  \item smaller transforms of length 7 and length 8
  \end{itemize}

  \medskip
  Resulting code is
  \begin{itemize}
  \item much faster: FFTs are $O(N \log N)$
  \item more accurate
  \end{itemize}

  \medskip
  Different FFT algorithms exist
  \begin{itemize}
  \item Cooley-Tukey
  \item Stockham
  \item Bluestein
  \end{itemize}
\end{frame}

\begin{frame}
  \frametitle{rocFFT}

  rocFFT highlights
  \begin{itemize}
  \item XXX
  \end{itemize}
\end{frame}

\begin{frame}
  \frametitle{rocFFT features vs the competition}

  XXX
\end{frame}

\section{Performance numbers}

\begin{frame}[label=perf01]
  \frametitle{Performance numbers}

  Important use cases (based on frequency of tickets)
  \begin{itemize}
  \item Molecular dynamics
  \item XXX Cholla
  \end{itemize}

  \medskip
  Performance tracking strategy
  \begin{itemize}
  \item Identify sizes of interest (from JIRA)
  \item Run performance suites; see \href{https://github.com/ROCmSoftwarePlatform/rocFFT-misc/tree/master/perf}{rocFFT-misc/perf}
  \item Identify and study performance gaps
  \end{itemize}
\end{frame}

\begin{frame}
  \frametitle{Performance numbers}

  2021-06-29 performance numbers

  \medskip
  \begin{tabular}{lc}
    one batch per transform: & \href{one-2021-06-29.html}{\beamerbutton{one}} \\
    multiple batches per transform: & \href{batched-2021-06-29.html}{\beamerbutton{batched}} \\
  \end{tabular}
\end{frame}

\section{Performance planning}

\begin{frame}
  \frametitle{Performance planning}

  To close performance gaps
  \begin{description}
  \item[FUSE] For multidimenional transforms, we are fusing transpose
    kernels with transform kernels.\par

    This is already mostly implemented.  Needs tuning for client
    transforms. Usually see 1.2x to 1.3x speedup.

  \item[NEWGEN] This is implemented.  Needs tuning for client
    transforms.  Often see 1.2x speedup; but can be higher.
  \end{description}

\end{frame}

\end{document}
